\section{Extended Related Work}
\label{app:related}
\label{sec:related}
% ============================================================================

\subsection{Chain-of-thought and its structural variants}

Chain-of-thought prompting~\citep{wei2022cot} and its zero-shot
form~\citep{kojima2022zeroshot} established that eliciting intermediate steps
improves reasoning accuracy, and self-consistency~\citep{wang2023selfconsistency}
established that sampling several chains and taking a plurality vote improves
it further. A family of structural elaborations followed:
least-to-most decomposition~\citep{zhou2023leasttomost}, tree-~\citep{yao2023tot}
and graph-structured~\citep{besta2024got} search over partial thoughts, and
program-mediated reasoning that offloads arithmetic to an
interpreter~\citep{chen2023pot,gao2023pal}. Later work tempered the picture:
\citet{sprague2025tocot} find that CoT's benefit is concentrated on
mathematical and symbolic tasks and is small or negative elsewhere, and much of
the apparent advantage of elaborate search structures narrows once sampling
budget is matched. \textbf{What remains unresolved:} the community still lacks
a principled account of \emph{which} items benefit from which structure, and
almost all such accounts are stated in terms of a per-item difficulty that, we
argue here, is itself unstable.

\subsection{Test-time compute scaling and its ceilings}

Our benchmarks are \textbf{MATH-500}~\citep{math500} (headline) and
GSM-Symbolic~\citep{mirzadeh2025gsmsymbolic}, with GSM8K~\citep{cobbe2021gsm8k}
as a secondary replicate. Verifier-reranked sampling~\citep{cobbe2021gsm8k}, process
supervision~\citep{lightman2024verify,uesato2022process,wang2024mathshepherd},
and compute-optimal allocation between sequential revision and parallel
sampling~\citep{snell2024scaling,wu2024inference} form the backbone of
test-time scaling; \citet{brown2024monkeys} showed that coverage---the
probability that at least one of $n$ samples is correct---keeps rising over
orders of magnitude of $n$, while plurality voting and learned reward models
plateau. \citet{muennighoff2025s1} showed that a very simple budget-forcing
scheme captures much of the benefit. Surveys now organise the
area along what/how/where/how-well axes~\citep{zhang2025ttssurvey} and by
subproblem topology~\citep{subproblemsurvey2025}, and a large matched-condition
study finds that no single strategy dominates across models and difficulty
regimes~\citep{artofscaling2025}.

The most directly relevant recent development is the explanation of the
plateau. \citet{brown2024monkeys} showed coverage keeps rising while plurality
voting plateaus; subsequent work formalises a \emph{modal ceiling}:
plurality voting converges to the mode, so selection accuracy saturates at
$\pmode$; the gap to rising coverage is the \emph{identifiability gap}
\citep{blendasc2025,certifiedsc2025,samplecomplexity2025}.
\textbf{What remains unresolved:}
all of this conditions on a per-item answer distribution under a \emph{fixed}
elicitation setting. \citet{bay2026modal} estimate run-to-run variance within a
configuration ($\rho_w \approx 0$) but never vary the prompt; whether $\pmode$
is configuration-relative is the step we test.

\subsection{Prompt sensitivity and variance decomposition}
\label{sec:related-prompt}

Three 2026 lines establish that per-item outcomes are prompt-attributable and
that evaluation pipelines must model crossed facets. \textbf{BrittleBench}
\citep{brittlebench2026} decomposes binary correctness variance into task
difficulty and semantics-preserving prompt sensitivity on an item\,$\times$\,perturbation
matrix; inference is deterministic (log-prob scoring), so the inference-variance
term vanishes. \textbf{TEE} \citep{tee2026} fits crossed random effects---including
item\,$\times$\,prompt---by REML and uses decision studies on aggregate benchmark
scores. \textbf{PromptEval} \citep{prompteval2024} estimates performance
distributions across prompt templates with an IRT-style model on the joint
(item, template) matrix. Earlier aggregate work
\citep{sclar2024formatspread,mizrahi2024multiprompt,alzahrani2024leaderboard}
showed format and template choice move rankings by large margins.

\textbf{What remains unresolved} is the intersection with \emph{sampled}
chain-of-thought on free-generation math and with the \emph{mode}. BrittleBench
does not report $\pmode$, reorder rates, or CDV; TEE's estimand is
$\mathrm{Var}(\hat\theta)$ at benchmark scale---item\,$\times$\,prompt is a
nuisance that averages away in aggregate CIs, whereas adaptive test-time compute,
hard-subset curation, and IRT consumers use \emph{per-item} parameters one at a
time. \citet{facetnoise2026,cyclicjudge2026} show that under G-theory facet decompositions,
repeats are the cheap facet and paraphrases/models buy more reliability per
query---the allocation-theoretic sibling of ``i.i.d.\ resampling cannot move the
plateau,'' but on brand sentiment without modes or reasoning. \citet{hardunreached2026}
establish unstable hardness labels along the stochastic-vs-deterministic axis;
we study the \emph{configuration} axis and whether it reaches the selection
bottleneck. Adjacent aggregation devices---paraphrase
voting~\citep{paraconsistency2024,zhou2024paraphrase}---are baselines, not the
measurement target.

\subsection{Long chain-of-thought, RL for reasoning, and the elicitation debate}

Reasoning-specialised models trained with reinforcement learning on verifiable
rewards~\citep{deepseekr1_2025,shao2024deepseekmath} produce long
chains-of-thought and large gains, and distillation transfers much of the
behaviour into small models with remarkably little
data~\citep{muennighoff2025s1,ye2025limo}. A pointed debate followed about what
RL actually contributes. \citet{yue2025rlvr} argue via pass@$k$ analysis that
RLVR largely sharpens the base model's existing distribution rather than
extending its reach, and \citet{diversitycollapse2025} formalise a related
bias--variance decomposition of pass@$k$, showing that supervised finetuning
improves pass@1 while collapsing generation diversity and degrading pass@$k$,
and that weight interpolation partially recovers both. Theoretically,
self-improvement has been framed as \emph{sharpening} a fixed base
policy~\citep{huang2025sharpening}, and inference-time methods have been
recast as sharpening the answer marginal
directly~\citep{marginalsharpening2026}. \textbf{What remains unresolved:} the
sharpening framing implies that the ordering of candidate answers is doing all
the work, but the stability of that ordering---under scale, under
post-training, and under elicitation---has not been measured. Our
mode-reordering analysis (\S\ref{sec:method-reorder}) is the missing
measurement for the elicitation axis.

\subsection{Faithfulness, premature commitment, and what the chain is for}

A parallel literature asks whether the chain reflects the computation.
\citet{turpin2023unfaithful} show that models silently follow biasing features
in the context while producing unrelated rationales;
\citet{lanham2023measuring} introduce the truncation and corruption probes that
have become standard; and \citet{lyu2023faithful} propose constrained
formulations that are faithful by construction. Latent-space variants ask
whether the tokens are needed at
all~\citep{hao2024coconut,pfau2024filler,goyal2024pause}. Recent work sharpens
the commitment question: \citet{prematureconfidence2026} probe at evenly spaced
truncation points, define a premature-confidence score from the resulting
confidence trajectory, show it predicts logical flaws, and mitigate it with a
shaped RL objective; \citet{precotprobe2026} show that a linear probe on
pre-CoT activations predicts the final answer with high AUC and that the
direction is causal under steering. \citet{matcha2025} show the converse
decoupling---answers that stay correct while the rationale collapses under
imperceptible input perturbation. \textbf{What remains unresolved:} these
results establish that a model often has an answer in mind early, but they do
not connect that to the \emph{selection} bottleneck at the heart of test-time
scaling. If the answer is fixed early, then what a configuration change
perturbs is that early commitment, which is precisely why it can move a mode
that resampling cannot.

\subsection{Efficiency, overthinking, and adaptive compute}

Reasoning models overspend on easy items~\citep{overthinking2024} and
under-explore on hard ones~\citep{underthinking2025}; surveys catalogue the
resulting efficiency
literature~\citep{stopoverthinking2025}. Length-controlled policy
optimisation~\citep{l1_2025} and budget-forcing~\citep{muennighoff2025s1} give
explicit control over chain length. On the parallel axis,
adaptive-consistency~\citep{aggarwal2023adaptive} and early-stopping
self-consistency~\citep{li2024esc} halt when the vote is decided;
Certaindex/Dynasor~\citep{fu2024certaindex} generalise this to an
algorithm-agnostic ``answer stabilisation'' signal used for scheduling;
Blend-ASC~\citep{blendasc2025} allocates samples across questions to optimise
sample efficiency; \citet{certifiedsc2025} provide sequential certificates; and
\citet{adaptivettc2026} study difficulty-conditioned budget allocation directly.
\textbf{What remains unresolved:} every one of these estimates a per-item
quantity in one setting and consumes it in another, and none validates that the
estimate transfers. This is the applied consequence we measure.

\subsection{Robustness, perturbation, and the reasoning-versus-memorisation debate}

\citet{mirzadeh2025gsmsymbolic} introduced templated variants of GSM8K and
reported both high variance across instantiations of the same template and a
large drop when an irrelevant clause is added (GSM-NoOp), arguing for
pattern-matching over formal reasoning. Contamination
audits~\citep{zhang2024gsm1k}, counterfactual
evaluation~\citep{wu2024reciting}, premise-order
effects~\citep{chen2024premise} and the reversal
curse~\citep{berglund2024reversal} point the same direction, while
compositionality analyses~\citep{dziri2023faithfate} give a mechanistic account
of multi-step failure. The debate is live: puzzle-complexity critiques and
their rebuttals continue, and 2026 replications of GSM-NoOp on frontier models
suggest that a substantial part of the original effect reflects models
reasonably treating ambiguous clauses as informative rather than failing to
reason. Follow-up
work finds that program-mediated reasoning does not repair GSM-Symbolic
robustness~\citep{reasoningcodeboth2026}. \textbf{What remains unresolved:}
these studies compare accuracy across perturbations of the \emph{problem}. Ours
holds the problem fixed and perturbs only the \emph{elicitation}, which
isolates a different and, we argue, more troubling source of instability,
because it cannot be blamed on the item.

\subsection{Self-correction, verification, and aggregation beyond frequency}

Iterative refinement with self-generated feedback~\citep{madaan2023selfrefine}
and multi-agent debate~\citep{du2024debate} were early attempts to convert extra
inference into extra accuracy without a verifier, but intrinsic self-correction
without external feedback does not reliably improve
reasoning~\citep{huang2024cannotselfcorrect,kamoi2024whencan}, which is part of
why the field turned to verifiers and to better aggregation. Beyond plain
plurality voting there is now a dense cluster: universal
self-consistency~\citep{chen2023universal}, program-verified
voting~\citep{prove2024}, semantic weighting of reasoning
paths~\citep{weightedreasoning2024}, confidence-informed voting and its
efficient variant~\citep{veccisc2026}, learned rerankers over answer-level
features~\citep{risc2026}, radial consensus scoring for
best-of-$N$~\citep{rcs2026}, and best-of-majority
hybrids~\citep{bestofmajority2025}.
Diversity-side interventions are covered in
\S\ref{sec:related-prompt}. \textbf{What remains unresolved:}
this literature searches for a better \emph{selector} given a fixed sample
pool. We ask a prior question---what determines the pool's mode in the first
place---and show that the answer includes a factor the selector literature
treats as fixed. The strongest of these methods are our baselines.

\subsection{Uncertainty, calibration, and evaluation methodology}

Confidence signals for reasoning include self-evaluated
correctness~\citep{kadavath2022know}, semantic
entropy~\citep{kuhn2023semantic,farquhar2024semantic}, verbalised
confidence~\citep{lin2022teaching,tian2023just}, and vote share itself. On the
evaluation side, the field has become considerably more careful:
\citet{miller2024errorbars} argues for routine interval reporting,
\citet{hochlehnert2025sober} document how much of recent reported progress on
small reasoning benchmarks is seed variance, \citet{dontpassk2025} propose
Bayesian posterior reporting in place of pass@$k$,
\citet{coverattau2025} show that pass@$k$ is biased toward low-reliability
``lucky hits'' and propose a reliability-controlled alternative, and
\citet{falsepositives2025} show that final-answer matching admits false
positives at a rate that materially lowers the true inference-scaling ceiling.
Psychometric approaches fit item difficulty explicitly and use it for adaptive
testing and routing~\citep{atlas2025,irtnet2025,irtrouter2025,psnirt2026}, and
\citet{kim2025correlated} show that errors are strongly correlated across
models, capping ensembles---an effect formalised as a co-failure
ceiling~\citep{cofailure2026}. \textbf{What remains unresolved:} the
psychometric line assumes a unidimensional item difficulty and validates it
largely on non-reasoning benchmarks. We test that assumption where it is most
consequential, and we adopt the evaluation hygiene this literature demands
throughout.

% ============================================================================
